Precise cosmological constraints from \textit{Planck} observations of Cosmic Microwave Background (CMB) anisotropies \citep{2020A&A...641A...6P} support a spatially flat \(\Lambda\)CDM Universe, which has proven remarkably successful in predicting a wide range of cosmological observables. These include the expansion history of the Universe, the growth of large-scale structure, and the matter distribution as inferred from weak gravitational lensing distortions of distant galaxies \citep{1972gcpa.book.....W, 1980lssu.book.....P, 1990ApJ...349L...1T, 1998ApJ...498...26K}.

Over the past decade, rapid progress in wide-field optical imaging surveys---notably the Kilo-Degree Survey \citep[KiDS;][]{2019A&A...625A...2K, 2024A&A...686A.170W}, the Dark Energy Survey \citep[DES;][]{2025arXiv250105665Y, 2025arXiv250105739B}, and the Hyper Suprime-Cam Survey \citep[HSC;][]{2022PASJ...74..247A, 2022PASJ...74..421L}---has substantially improved late-Universe tests of this paradigm. These surveys measure the auto-correlations of background galaxy ellipticities \citep{2021A&A...645A.104A, 2022PhRvD.105b3514A, 2022PhRvD.105b3515S, 2023A&A...679A.133L, 2023PhRvD.108l3519D, 2023PhRvD.108l3518L}, commonly referred to as cosmic shear \citep{2001PhR...340..291B, 2015RPPh...78h6901K}. Their constraining power is further enhanced by cross-correlations with foreground galaxy positions (galaxy--galaxy lensing) and by measurements of foreground galaxy clustering. Together, these three two-point statistics form the basis of the so-called \(\mathit{3 \times 2 \, pt}\) analysis \citep{2021A&A...646A.140H, 2022PhRvD.105b3520A, 2023A&A...675A.189D, 2023PhRvD.108l3521S, 2023PhRvD.108l3517M}.

Collectively, these probes have delivered progressively tighter constraints on the structure-growth parameter \(S_8 = \sigma_8 \sqrt{\Omega_\text{m}/0.3}\), providing an important and independent test of the standard \(\Lambda\)CDM framework. Earlier analyses suggested a \(\sim2\)--\(3\sigma\) discrepancy---its precise level depending on the combination of surveys and modelling choices---between late-Universe measurements and early-Universe constraints from \textit{Planck} \citep{2023OJAp....6E..36D, 2023MNRAS.518..477A, 2023MNRAS.520.5016L, 2023A&A...675A.189D, 2023PhRvD.108l3519D, 2024JCAP...08..024G}. More recently, however, the KiDS-Legacy cosmic shear \citep{2025arXiv250319442S, 2025arXiv250319441W} and DES-Y6 analyses \citep{2026arXiv260210065D, 2026arXiv260114559D} have reported values in statistically plausible agreement with \textit{Planck}. This improved concordance largely results from increasingly sophisticated modelling of systematic effects, especially advances in spectroscopic coverage and photometric-redshift calibration. This trajectory underscores the central role of rigorous systematic control as a precondition for unlocking the cosmological information of next-generation surveys.

The coming decade will mark a transformative era for wide-field imaging cosmology. The NSF--DOE Vera C. Rubin Observatory \citep[\textit{Rubin};][]{2009arXiv0912.0201L}, together with space-based missions such as \textit{Euclid} \citep{2025A&A...697A...1E}, the Nancy Grace Roman Space Telescope \citep{2015arXiv150303757S}, and the China Space Station Telescope (CSST) \citep{2019ApJ...883..203G}, will provide unprecedented imaging depth, area, and image quality. These surveys will observe billions of galaxies, rendering complete spectroscopic follow-up infeasible. As a scalable alternative, cosmological analyses rely on photometric-redshift techniques, which infer redshifts from broadband flux measurements \citep{1985AJ.....90..418K, 1986ApJ...303..154L}. Achieving the full statistical power of these surveys requires exceptional control of systematic uncertainties \citep{2018ARA&A..56..393M}, among which photometric-redshift estimation and calibration remain particularly demanding.

As reviewed by \citet{2022ARA&A..60..363N}, photometric redshifts enter \(3\times2\,\text{pt}\) analyses in two distinct ways. First, individual galaxy point estimates (e.g.\ maximum-likelihood or posterior-mean redshifts) are used to define foreground lens and background source samples, and to assign galaxies to tomographic redshift bins \citep{1999ApJ...522L..21H, 2002PhRvD..65f3001H}. The fidelity of these point estimates governs the level of contamination from catastrophic outliers and misclassified galaxies, which can propagate into biases in shear and clustering correlations \citep{2006ApJ...636...21M}. Secondly, and more directly relevant to cosmological inference, the population-level redshift distribution associated with a given tomographic bin enters the theoretical modelling of both lensing and clustering observables. Even small biases in its mean or shape can induce significant shifts in cosmological parameters \citep{2012MNRAS.421.2355H, 2013MNRAS.431.1547B}. Consequently, while point estimates determine the purity of tomographic binning, it is the accuracy of tomographic redshift distributions that most directly controls cosmological inference. Both aspects are explicitly recognised in the \textit{Rubin} Legacy Survey of Space and Time (LSST) Dark Energy Science Collaboration (DESC) Science Requirements Document \citep[SRD;][]{2018arXiv180901669T}, which sets stringent performance targets for redshift estimation and calibration.

Substantial effort has been devoted to improving point-estimate photometric redshifts. Broadly, methods can be divided into template-fitting approaches, empirical machine-learning (ML) techniques, and hybrid schemes. Template-fitting methods match observed fluxes to model spectral energy distributions (SEDs) \citep[e.g.][]{1999MNRAS.310..540A, 2000ApJ...536..571B, 2006A&A...457..841I, 2006AJ....132..926C, 2006MNRAS.372..565F, 2008ApJ...686.1503B, 2017ApJ...838....5L}, offering physical interpretability and potential extrapolation capability. ML approaches learn mappings from photometric features to redshift using spectroscopic training data \citep[e.g.][]{2004PASP..116..345C, 2012A&A...546A..13C, 2016PASP..128j4502S, 2016MNRAS.462..726A, 2017MNRAS.465.1959C, 2018MNRAS.473.3969R, 2020MNRAS.499.1587S, 2024AJ....168...80C, 2025OJAp....8E..50M}. Hybrid methods aim to combine the complementary strengths of these strategies \citep[e.g.][]{2019ApJ...881...80L, 2024MNRAS.530.2012T, 2025A&A...698A.276J}. Systematic benchmarks of template-fitting, ML, and hybrid algorithms against common datasets have been carried out in preparation for Stage-IV surveys \citep{2020A&A...644A..31E}.

The most widely adopted approach for calibrating tomographic redshift distributions is the Direct Calibration (\texttt{DIR}) technique \citep{2008MNRAS.390..118L, 2009MNRAS.396.2379C}, later refined with \(k\)-Nearest Neighbour methods \citep{2016MNRAS.463..635H, 2020A&A...633A..69H}. SOM-based partitioning of colour space offers a complementary clustering approach \citep{2020A&A...637A.100W, 2025arXiv250319440W, 2026A&A...707A.277R}. Multi-layer SOM frameworks further exploit deep-drilling fields (DDFs)---small-area pointings that reach several magnitudes deeper than the main survey and benefit from more complete spectroscopy---to anchor the shallow-sample calibration through a transfer relation established between the deep and wide photometry \citep{2019MNRAS.489..820B, 2021MNRAS.505.4249M, 2024arXiv240800922C, 2025arXiv250907964G, 2025arXiv251023566Y}. Alternative strategies include stratified learning approaches \citep{2024MNRAS.534.3808A} and clustering-redshift techniques based on angular cross-correlations \citep{2008ApJ...684...88N, 2020A&A...642A.200V, 2023MNRAS.524.5109R, 2025A&A...702A.155D}. Nevertheless, all such approaches remain fundamentally limited by the completeness and representativeness of spectroscopic calibration data. Gaps in colour--magnitude space, particularly for faint or high-redshift galaxies, can propagate into biases in the inferred tomographic redshift distributions \citep{2015ApJ...813...53M, 2020MNRAS.496.4769H}, motivating simulation-based strategies that augment the spectroscopic calibration sample.

To address this incompleteness, \citet{2024ApJ...967L...6M} proposed augmenting spectroscopic training sets with simulated galaxies drawn from mock catalogues, providing the first proof of concept that such data augmentation can improve photometric-redshift performance for LSST-like surveys. Building on this proposal, in our previous work (hereafter Paper~I; \citealt{2026MNRAS.547f2226Z}), we introduced a self-organising map (SOM)-based augmentation framework that uses the SOM manifold to systematically identify regions of colour space where spectroscopic information is missing, enabling a more flexible and targeted augmentation strategy. Paper~I further provided a more realistic validation of the augmentation approach under LSST Year~1 and Year~10 observing scenarios, benchmarking point-estimate photometric-redshift performance against LSST science requirements. 

In the present work, we extend this SOM-based augmentation framework from galaxy-level redshift estimation to the calibration of population-level tomographic redshift distributions. By incorporating synthetic galaxy SEDs into the SOM framework, we improve the representativeness of the spectroscopic calibration data and derive conditional redshift distributions (here, the condition refers to the selection function defining each tomographic bin; see Section~\ref{sec:2} for the formal definition) via importance weighting within grouped clusters of SOM cells. We combine these direct estimates with ML-based density modelling to construct hybrid estimators. Furthermore, we introduce a probabilistic sampling framework that explicitly incorporates spectroscopic selection functions, enabling forward propagation of both statistical and systematic uncertainties into experiment-marginalised tomographic redshift distributions for lens and source samples. Finally, we derive prior information for commonly adopted parameterisations---such as mean-redshift shifts and width rescaling---and discuss the limitations of our current analysis and the implications for future work.

This paper is organised as follows. In Section~\ref{sec:2}, we establish the Bayesian framework used throughout and derive practical estimators for conditional redshift distributions. Section~\ref{sec:3} then specifies the experimental design, including the construction of the photometric and spectroscopic datasets, the tomographic selections, and the set of alternative calibration configurations used to probe different assumptions about support, augmentation, and selection. In Section~\ref{sec:4}, we compare the resulting estimators and assess their calibration performance across the ensemble of realised experiments. Building on this, Section~\ref{sec:5} changes perspective: rather than treating each experiment separately, it uses the full ensemble of realisations to define a probabilistic uncertainty framework, construct experiment-marginalised tomographic redshift distributions, and derive analysis-ready uncertainty products. Section~\ref{sec:6} discusses the main limitations of the framework and their implications for redshift calibration in Stage-IV weak-lensing surveys. Finally, Section~\ref{sec:7} summarises our main findings and conclusions.